\section{Positioning in the literature}\label{sec:lit-table}

Table~\ref{tab:littable} summarizes how the case-study result relates to the literatures reviewed in Section~\ref{sec:lit} of the paper.


\begin{table}[ht]
\centering\small
\begin{tabularx}{\textwidth}{@{}>{\raggedright\arraybackslash}p{3.1cm}>{\raggedright\arraybackslash}X>{\raggedright\arraybackslash}X@{}}
\toprule
\textbf{Literature} & \textbf{What it characterizes} & \textbf{Difference from this paper}\\
\midrule
Yaari (1987); Wang, Young, Panjer (1997)
 & One premium functional $H$; axioms on $H$ (law invariance, monotonicity, comonotonic additivity).
 & We characterize a \emph{ratio} of two such functionals with axioms on the ratio. Monotonicity is placed on the \emph{implied premium and risk measure}, never on the ratio.\\
Schmeidler (1986); Murofushi--Sugeno--Machida (1994); Cerreia-Vioglio et al.\ (2012, 2015); Wang--Wei--Willmot (2020)
 & Representation of one comonotonic additive functional as a (signed) Choquet integral.
 & Used as the engine (applied to $\rho$ and $\Pi$). New: the axioms are on $\LR$, comonotonic additivity of $\rho$ and $\Pi$ is \emph{derived} (also on the break-even set), and uniqueness is unconditional.\\
Cherny--Madan (2009); Bielecki--Cialenco--Zhang (2014); Biagini--Bion-Nadal (2014)
 & Monotone, quasi-concave indices $\Leftrightarrow$ increasing families of coherent measures.
 & $\LR$ cannot be monotone (Prop.~\ref{prop:F2}); the associated family is signed and non-monotone. Cardinal (pooling identity) rather than ordinal.\\
Cheridito--Kromer (2013)
 & Axioms for reward--risk ratios (sufficient conditions).
 & Reward over risk is monotone; loss over premium is not. We give an \emph{iff}, within a stated class.\\
Han--Lin--Wang (2025)
 & Diversification quotient: iff axiomatization of a level-crossing index.
 & Closest in spirit; monotone families and a portfolio index. We characterize the ratio and its parts.\\
Haezendonck--Goovaerts (1982); Li--Wang (2023); Bellini et al.\ (2018)
 & Implicitly defined premia/risk measures, return risk measures.
 & Existence and uniqueness of an implicit functional, not a ratio representation. The association is recovered as a corollary and shown to hold only off $\{\rho=\Pi=0\}$.\\
Chew (1983); Dekel (1986)
 & Ratio (weighted utility) versus implicit (betweenness) representations.
 & Explains why a mediant inequality is not enough and a premium-weighted identity is needed.\\
Gneiting (2011)
 & Elicitability of ratios of expectations.
 & Statistical, no axioms on pooling or cash shifts.\\
\bottomrule
\end{tabularx}
\caption{Positioning of the case-study result. ``Iff'' means a characterization by necessary and sufficient axioms.}\label{tab:littable}
\end{table}